\section{Signalling dynamics: what live-cell traces encode}
\label{sec:partB}
The Tay-lab deposit \citep{wang2025} is the one public resource in which live-cell
dynamics of macrophages, here single-cell NF-\(\kappa\)B reporter traces, are paired at the
condition level with transcriptome and chromatin data, and it includes sequential
stimulations, which is the structure innate memory has.

\paragraph{Ligand-pair identity is decodable from dynamics.} Across \val{pb.lig.wells}
wells and \val{pb.lig.n} cells, the identity of the ordered pair of stimuli
(\val{pb.lig.classes} classes, chance \val{pb.lig.chance}) is decoded from single-cell
traces well above chance, with wells as groups (Figure~\ref{fig:signalling}): the raw trace
reaches \val{pb.lig.raw_trace} (95\% CI \val{pb.lig.raw_trace.ci}), hand-crafted trace
features \val{pb.lig.handcrafted} (\val{pb.lig.handcrafted.ci}), and a trace JEPA, pretrained
without labels on \val{pb.pre.cells} traces from \val{pb.pre.wells} wells that are not
evaluated, \val{pb.lig.trace_jepa} (\val{pb.lig.trace_jepa.ci}). The pretraining is not
what helps: the same network at random initialisation reaches
\val{pb.lig.trace_jepa_random_init} (difference \val{pb.lig.d.trace_jepa_random_init},
95\% CI \val{pb.lig.d.trace_jepa_random_init.ci}), and its advantage over the raw trace
(\val{pb.lig.d.raw_trace}, \val{pb.lig.d.raw_trace.ci}) does not survive correction for
multiple comparisons. Hand-crafted summaries lose information that the full trace carries
(\val{pb.lig.d.handcrafted}, \val{pb.lig.d.handcrafted.ci}).

\paragraph{Acute tolerance cannot be tested with the deposited wells.} The deposit has two
wells of LPS-then-LPS and two of naive-then-LPS cells, so a well-level test of tolerance
from dynamics has no power. From the observed well means of the second response's peak
amplitude (standardised effect \val{pb.pow.peak_amplitude.d}) and their between-well SD, a
two-sided test at $\alpha=0.05$ needs \val{pb.pow.peak_amplitude.n} wells per arm for 80\%
power, and \val{pb.pow.peak_amplitude.nhi} if the SD is at the upper limit of its 80\%
interval; for the response's area under the curve the numbers are \val{pb.pow.auc.n} and
\val{pb.pow.auc.nhi}. Because that SD rests on four wells, we read this as a lower bound on
what a tolerance-from-dynamics study needs, not as a precise design.

\paragraph{Dynamics to omics and across species.} With five matched conditions, condition-
level dynamics do not predict held-out transcript or accessibility profiles better than the
leave-one-out mean (Table~\ref{tab:coverage}), and among \val{b4.both.uni} orthologs induced
by LPS in both mouse macrophages and human macrophages, the genes blunted by heterologous
priming in mouse overlap the genes blunted by homologous tolerance in human no more than
expected by chance (\val{b4.both.k} observed, \val{b4.both.exp} expected, odds ratio
\val{b4.both.or}, two-sided $p=\val{b4.both.p}$). Both analyses are exploratory.
